\section{Limiti e funtori}
\begin{proposition}[Limiti e colimiti in categorie di funtori]\label{prop_limiti_colimiti_categorie_funtori}\index{Limite!in categorie di funtori}\index{Colimite!in categorie di funtori}\index{Categoria!di funtori}
	\Todo{}
\end{proposition}

\begin{remark}\label{mk_colim_prefasci_pi0_elts}\index{Colimite!in \(\ctSet\)}\index{Prefascio!colimite}\index{\(\pi_0\)}\index{Elementi!categoria degli ---}
	Il caso particolare di limiti in \(\ctSet\) (prefasci): per \(F \colon \ctC^\op\fun\ctSet\), \(\colim F\cong \pi_0(\Elts\ctC F)\). E il limite?
\end{remark}

\begin{definition}[Preservazione, riflessione, creazione di limiti]\label{def_preservazione_riflessione_creazione}\index{Limite!preservazione}\index{Limite!riflessione}\index{Limite!creazione}\index{Colimite!preservazione}\index{Colimite!riflessione}\index{Colimite!creazione}
	\Todo{}
\end{definition}

\begin{proposition}[Il dimenticante da \(\ctAlg(F)\) crea i limiti]\label{prop_dimenticante_alg_crea_limiti}\index{Algebra!di un funtore}\index{Coalgebra!di un funtore}\index{Funtore!dimenticante!creazione di limiti}\index{Limite!creazione}\index{Colimite!creazione}
	\Todo{}
\end{proposition}

\begin{theorem}\label{thm_yoneda_preserva_limiti}\index{Yoneda!immersione di ---!preserva i limiti}\index{Limite!preservato da Yoneda}
	\Todo{Yoneda preserva i limiti}
\end{theorem}
\begin{theorem}
	\Todo{Funtori pienamente fedeli e riflessione di co/limiti}
\end{theorem}
\begin{theorem}
	\Todo{Funtore dimenticante che crea sia limiti che colimiti}
\end{theorem}
\begin{corollary}
	\Todo{Co/limiti in categorie di funtori	si calcolano punto a punto (già fatto sopra, ora ottenibile in modo elegante).}
\end{corollary}
\begin{definition}[Teorie-prodotto]\label{def_teorie_prodotto}\index{Teoria!algebrica}\index{Teoria!prodotto}
	% i.e. teorie algebriche
	\Todo{}
\end{definition}

\begin{definition}[Teorie-finlimite]\label{def_teorie_finlimite}\index{Teoria!essenzialmente algebrica}\index{Teoria!finlimite}
	% i.e. teorie essenzialmente algebriche
	\Todo{}
\end{definition}

\begin{theorem}[Relazione di Yoneda con le teorie limite]\label{thm_yoneda_teorie_limite}\index{Yoneda!teorie limite}\index{Teoria!limite}
	\Todo{}
\end{theorem}

\subsection{Coincidenza di limiti e colimiti}

\begin{definition}[Biprodotti finiti in \(\ctAb\)]\label{def_biprodotti}\index{Biprodotto}
	\Todo{}
\end{definition}
Questa proprietà è condivisa da molte altre categorie oltre a \(\ctAb\); più in generale, tutte le categorie \(\ctMod(R)\) di moduli su un anello hanno biprodotti, e più in generale ancora, tutte le categorie abeliane (si veda \ref{def_categoria_abeliana}) hanno biprodotti. In effetti, la presenza di biprodotti è caratterizzante delle categorie abeliane, che sono particolari categorie \emph{additive} (si veda \ref{def_categoria_additiva}).
\begin{definition}[Quadrati esatti]\label{def_quadrati_esatti}\index{Quadrato esatto}\index{Quadrato esatto!in \(\ctAb\)}
	\Todo{}
\end{definition}
\begin{proposition}
	\Todo{Limiti assoluti sono anche colimiti, or something}
\end{proposition}
\begin{proposition}[Esattezza e successioni esatte corte]\label{prop_esattezza_successioni_esatte_corte}\index{Successione esatta corta}\index{Quadrato esatto!relazione con successioni esatte}
	\Todo{}
\end{proposition}

\begin{definition}[Categoria abeliana]\label{def_categoria_abeliana}\index{Categoria!abeliana}
	\Todo{}
\end{definition}

\subsection{Proprietà di oggetti specificate da limiti}
\begin{definition}[Diagramma iniziale, diagramma finale]\label{def_diagramma_iniziale_finale}\index{Diagramma!iniziale}\index{Diagramma!finale}\index{Funtore!iniziale}\index{Funtore!finale}
	\Todo{}
\end{definition}
(il contenuto della sezione \ref{} su funtori iniziali e finali ora è in contesto)

\begin{definition}[Funtore denso, sottocategoria densa, funtore codenso, sottocategoria codensa]\label{def_funtore_denso_codenso}\index{Funtore!denso}\index{Funtore!codenso}\index{Sottocategoria!densa}\index{Sottocategoria!codensa}
	\Todo{}
\end{definition}

\begin{definition}[Oggetto \(\kappa\)-compatto]\label{def_oggetto_kappa_compatto}\index{Oggetto!\(\kappa\)-compatto}\index{Compattezza}
	\Todo{}
\end{definition}

\begin{definition}[Oggetto minuscolo]\label{def_oggetto_minuscolo}\index{Oggetto!minuscolo}
	\Todo{}
\end{definition}

\begin{definition}[Oggetto connesso]\label{def_oggetto_connesso}\index{Oggetto!connesso}
	\Todo{}
\end{definition}

\begin{definition}[Oggetto \(\bbD\)-presentabile]\label{def_oggetto_D_presentabile}\index{Oggetto!\(\bbD\)-presentabile}\index{Presentabilità}
	\Todo{}
\end{definition}